\documentclass[../main.tex]{subfiles}
\ifSubfilesClassLoaded{\setcounter{chapter}{3}}{}
\begin{document}

\chapter{Struktur Kontrol Percabangan}

\begin{subcpmk}
  \item Sub-CPMK 2.1: Menyusun program dengan \texttt{if}/\texttt{elif}/\texttt{else} untuk percabangan
\end{subcpmk}

\SesiRPS{4}{2.1}{$\sim$170 menit}{CPMK-2}

\section{Sarana dan prasarana}

Python 3.10+, editor, terminal.

\section{Prosedur kerja}

\begin{enumerate}[leftmargin=*]
  \item Pelajari pola percabangan dan \texttt{match} pengenalan.
  \item Kerjakan soal bertahap dari sederhana ke bersarang.
  \item Verifikasi dengan beberapa kasus uji manual.
\end{enumerate}

\section{Sintaks Percabangan \texttt{if}/\texttt{elif}/\texttt{else}}

Struktur dasar percabangan di Python:

\begin{lstlisting}[caption=if / elif / else (Tutorial Python ss4.1)]
nilai = int(input("Nilai: "))
if nilai >= 80:
    print("A")
elif nilai >= 70:
    print("B")
elif nilai >= 60:
    print("C")
else:
    print("E")
\end{lstlisting}

Indentasi menandai blok yang termasuk dalam cabang \keyword{if}. Tidak ada tanda kurung kurawal seperti bahasa lain --- indentasi \textbf{adalah} struktur program.

\section{Kondisi Bertingkat dan Bersarang}

Percabangan dapat bersarang untuk logika yang lebih kompleks. Hindari terlalu dalam tanpa komentar; pertimbangkan memecah ke fungsi pada modul berikutnya.

\begin{lstlisting}[caption=Contoh percabangan bersarang]
x = int(input("x: "))
if x >= 0:
    if x == 0:
        print("nol")
    else:
        print("positif")
else:
    print("negatif")
\end{lstlisting}

\section{Ekspresi Kondisional (Ternary)}

Python mendukung bentuk singkat \texttt{a if kondisi else b} yang menghasilkan nilai berdasarkan kondisi:

\begin{lstlisting}[caption=Ekspresi ternary]
x = int(input("x: "))
label = "genap" if x % 2 == 0 else "ganjil"
print(x, "adalah", label)
\end{lstlisting}

\begin{catatan}
Prioritas logika harus jelas: gunakan tanda kurung pada kondisi kompleks agar mudah dibaca.
\end{catatan}

\section{Pernyataan \texttt{pass}}

Pernyataan \keyword{pass} tidak melakukan apa-apa. Berguna sebagai \textit{placeholder} ketika sintaks mensyaratkan sebuah pernyataan tetapi program belum memerlukan aksi apapun. Ini sering digunakan selama pembuatan kode (prototyping).

\begin{lstlisting}[caption=Penggunaan pass sebagai placeholder (Tutorial Python ss4.6)]
def fungsi_belum_jadi():
    pass   # akan diimplementasikan nanti

# Dalam blok if yang belum diisi:
pilihan = input("Pilih (a/b): ")
if pilihan == "a":
    print("Anda memilih a")
elif pilihan == "b":
    pass   # TODO: implementasikan opsi b
else:
    print("Pilihan tidak valid")
\end{lstlisting}

\section{Pernyataan \texttt{match} (Pattern Matching)}

Sejak Python 3.10, tersedia pernyataan \keyword{match} sebagai alternatif dari rantai \texttt{if}/\texttt{elif} yang panjang. \texttt{match} mencocokkan nilai dengan sejumlah \textit{pattern} (pola) dan membuat kode lebih mudah dibaca.

\begin{lstlisting}[caption=Pernyataan match dasar (Tutorial Python ss4.7)]
perintah = input("Perintah (berhenti/mulai/jeda): ")
match perintah:
    case "berhenti":
        print("Program dihentikan.")
    case "mulai":
        print("Program dimulai.")
    case "jeda":
        print("Program dijeda.")
    case _:              # wildcard, cocok dengan apapun
        print(f"Perintah '{perintah}' tidak dikenal.")
\end{lstlisting}

\subsection{match dengan Beberapa Pola dalam Satu Case}

Beberapa pola dapat digabung menggunakan \texttt{|} (OR):

\begin{lstlisting}[caption=Gabungan pola dengan OR]
kode_status = 404
match kode_status:
    case 200 | 201 | 204:
        print("Sukses")
    case 400:
        print("Permintaan tidak valid")
    case 404:
        print("Tidak ditemukan")
    case 500 | 503:
        print("Kesalahan server")
    case _:
        print("Kode tidak dikenal")
\end{lstlisting}

\begin{konsep}
\texttt{match} lebih ekspresif daripada rantai \texttt{if/elif} saat mencocokkan satu variabel dengan banyak nilai tetap. Namun \texttt{if/elif} tetap diperlukan untuk kondisi dengan ekspresi boolean yang lebih kompleks (misalnya rentang nilai).
\end{konsep}

\begin{aktivitas}
  \item Buat program yang menentukan bilangan ganjil/genap dari masukan integer.
  \item Buat menu sederhana (1: luas persegi, 2: luas lingkaran) menggunakan \texttt{if}/\texttt{elif}, kemudian ubah versinya dengan \texttt{match}.
  \item Implementasikan aturan nilai huruf (A--E) dari rentang nilai angka menggunakan \texttt{if}/\texttt{elif}.
  \item Buat program kalkulator sederhana yang membaca dua angka dan satu operator (\texttt{+}, \texttt{-}, \texttt{*}, \texttt{/}) lalu gunakan \texttt{match} untuk menentukan operasi yang dijalankan.
  \item Tambahkan blok \texttt{pass} sebagai placeholder untuk menu yang belum diimplementasikan.
\end{aktivitas}

\begin{latihan}
  \item Jelaskan kapan \texttt{elif} lebih tepat daripada banyak \texttt{if} independen.
  \item Apa perilaku program jika semua kondisi \texttt{if}/\texttt{elif} salah dan tidak ada \texttt{else}?
  \item Tulis pseudokode lalu kode Python untuk membandingkan tiga bilangan dan mencetak yang terbesar.
  \item Kapan lebih baik menggunakan \texttt{match} daripada \texttt{if}/\texttt{elif}? Sebutkan satu kasus yang cocok untuk masing-masing.
  \item Apa fungsi \texttt{case \_:} dalam blok \texttt{match}? Apa padanannya dalam struktur \texttt{if}/\texttt{elif}?
  \item Refleksi: bagaimana percabangan membantu menangani kasus ujung (\textit{edge cases})?
\end{latihan}

\begin{asesmen}
\textbf{Praktik di lab}

\begin{enumerate}
  \item Diberikan \texttt{x} bilangan bulat, kondisi manakah yang benar untuk ``x habis dibagi 3 tetapi bukan habis dibagi 2''?
  \begin{enumerate}
    \item \texttt{x \% 3 == 0 and x \% 2 != 0}
    \item \texttt{x \% 3 == 0 or x \% 2 != 0}
    \item \texttt{x // 3 == 0}
    \item \texttt{x \% 6 == 0}
  \end{enumerate}
  \item Pernyataan \texttt{pass} digunakan untuk:
  \begin{enumerate}
    \item Menghentikan loop
    \item Melewati iterasi berikutnya
    \item Menjadi placeholder kosong yang valid secara sintaks
    \item Mengembalikan nilai dari fungsi
  \end{enumerate}
  \item Fitur \texttt{match} tersedia mulai Python versi:
  \begin{enumerate}
    \item 3.8
    \item 3.9
    \item 3.10
    \item 3.12
  \end{enumerate}
\end{enumerate}

\textbf{Tugas}: program \texttt{tugas03.py} yang membaca tahun dan menentukan tahun kabisat (aturan: habis dibagi 4, kecuali abad kecuali habis dibagi 400). Kemudian buat program kalkulator \texttt{kalkulator.py} yang menggunakan \texttt{match} untuk memilih operasi aritmatika.
\end{asesmen}

\begin{checklist}
  \item Saya dapat menulis rantai \texttt{if}/\texttt{elif}/\texttt{else} yang benar secara indentasi
  \item Saya dapat menggabungkan kondisi dengan operator logika
  \item Saya dapat menggunakan \texttt{pass} sebagai placeholder blok kosong
  \item Saya dapat menuliskan \texttt{match} untuk pencocokan pola sederhana
  \item Saya dapat menggabungkan beberapa pola dengan \texttt{|} dalam \texttt{case}
  \item Saya dapat menyelesaikan masalah pemilihan dengan lebih dari dua cabang
\end{checklist}

\begin{rangkuman}
Percabangan mengarahkan eksekusi berdasarkan kondisi boolean. \texttt{elif} mengurangi \textit{nesting} yang tidak perlu. \texttt{pass} adalah pernyataan kosong yang berguna saat prototyping. \texttt{match} (Python 3.10+) menyederhanakan pencocokan nilai tunggal dengan banyak kasus.

\textbf{Referensi:} {\small Python Tutorial §4.1--4.7 \url{https://docs.python.org/3/tutorial/controlflow.html}}

\textbf{Kata kunci:} \code{if}, \code{elif}, \code{else}, \code{pass}, \code{match}, \code{case}, \code{and}, \code{or}, \code{not}
\end{rangkuman}

\section{Referensi sesi}

\begin{itemize}[leftmargin=*]
  \item \textbf{[5]} Univ. Pancasila, \textit{Modul praktikum Python} (percabangan).
  \item \textbf{[17]} Programiz, \textit{Python Programming Tutorial}.
  \item \textbf{[8, 9]} \textit{The Python Tutorial}, bagian alur kontrol.
\end{itemize}

\end{document}
